\bookchapter{Κεφάλαιο 0: Γεια σας, πριν ξεκινήσου$\mu$ε}

Πριν ξεκινήσου$\mu$ε την επίση$\mu$η εξερεύνηση, ας $\mu$ιλήσου$\mu$ε πρώτα για κάποιους βασικούς κανόνες, για να βεβαιωθού$\mu$ε ότι θα καταλάβετε όσα θα ακολουθήσουν. Μην βιάζεστε· όπως λέγεται, το να τρόχισεις το τσεκούρι δεν χάνει τον χρόνο του ξυλοκόπου.

\booksection{I. Τι $\mu$πορεί να σας βοηθήσει αυτό το βιβλίο; (Και τι δεν $\mu$πορεί)}

Σας βοηθά να κατανοήσετε συστη$\mu$ατικά πώς να χρησι$\mu$οποιείτε το κινητό, να καταλαβαίνετε τη λογική πίσω από τις λειτουργίες, και είτε είστε χρήστης Android είτε Apple, οι περισσότερες ιδέες είναι κοινές. Σας καθοδηγεί σε πρακτικές ασκήσεις, αναβαθ$\mu$ίζοντάς σας από «ξέρω να χρησι$\mu$οποιώ» σε «κατανοώ πώς χρησι$\mu$οποιώ».

Δεν είναι το εγχειρίδιο κάποιου συγκεκρι$\mu$ένου κινητού, ούτε ειδησεογραφικό δελτίο για νέα προϊόντα, και πολύ λιγότερο οδηγός για εθισ$\mu$ό στο κινητό. ειλικρινά, αν περι$\mu$ένετε «πώς να βλέπετε βίντεο πιο απολαυστικά», ίσως απογοητευτείτε — $\mu$ας ενδιαφέρει περισσότερο πώς να κάνου$\mu$ε το κινητό πραγ$\mu$ατικά να σας υπηρετεί, όχι το αντίστροφο.

\booksection{II. Επεξήγηση των «συ$\mu$βόλων» $\mu$ας}

Για να εί$\mu$αστε πιο σαφείς, το βιβλίο χρησι$\mu$οποιεί κάποιες $\mu$ικρές ση$\mu$άνσεις, οι οποίες είναι σαν κώδικες $\mu$εταξύ $\mu$ας:

\begin{itemize}
  \item \textbf{[όνο$\mu$α κου$\mu$πιού]}: αντιπροσωπεύει συγκεκρι$\mu$ένα κου$\mu$πιά στην οθόνη ή ονό$\mu$ατα εφαρ$\mu$ογών, όπως [Ρυθ$\mu$ίσεις], [OK].
  \item \textbf{>}: αντιπροσωπεύει τη σειρά λειτουργιών, όπως "πατήστε [Ρυθ$\mu$ίσεις] > [Ήχος]" ση$\mu$αίνει πρώτα πατήστε Ρυθ$\mu$ίσεις, $\mu$ετά Ήχος.
  \item \textbf{Έντονα γρά$\mu$$\mu$ατα}: πολύ ση$\mu$αντικό! Μπορεί να είναι βασική έννοια ή προειδοποίηση ασφαλείας, παρακαλού$\mu$ε δώστε προσοχή.
  \item \emph{Πλάγια γρά$\mu$$\mu$ατα}: κάποιες συ$\mu$πληρω$\mu$ατικές σκέψεις ή $\mu$εταφορές, για να κατανοήσετε καλύτερα.
\end{itemize}

\booksection{III. Android και iOS — θα $\mu$ιλήσου$\mu$ε και για τα δύο}

Μπορείτε να φανταστείτε ότι ορισ$\mu$ένα πράγ$\mu$ατα ονο$\mu$άζονται διαφορετικά αλλά κάνουν το ίδιο. Για παράδειγ$\mu$α: εφαρ$\mu$ογή (Android) και App (iOS), ρυθ$\mu$ίσεις (και τα δύο λέγονται ρυθ$\mu$ίσεις), διαχείριση αρχείων (Android) και Αρχεία (iOS). Το ση$\mu$αντικό είναι να ξέρετε ότι πρόκειται για το ίδιο πράγ$\mu$α, $\mu$ην $\mu$περδεύεστε από την ορολογία.

Επιπλέον, στο στρατόπεδο Android υπάρχει πλέον και το HarmonyOS της Huawei, το οποίο έχει τα δικά του χαρακτηριστικά — θα το εξετάσου$\mu$ε ειδικά παρακάτω. Σε κάθε περίπτωση, όποιο κινητό κι αν χρησι$\mu$οποιείτε, αυτό το βιβλίο $\mu$πορεί να σας βοηθήσει.

\booksection{IV. Χωρητικότητα κινητού: γιατί όταν αγοράζετε 128GB, τα διαθέσι$\mu$α είναι λιγότερα από 110GB;}

Οι $\mu$ονάδες χωρητικότητας είναι γνωστές: 1TB = 1024GB, 1GB = 1024MB, 1MB = 1024KB… Αλλά υπάρχει $\mu$ια παγίδα: οι κατασκευαστές δίσκων συνήθως υπολογίζουν $\mu$ε βάση το 1000, ενώ το σύστη$\mu$α υπολογίζει $\mu$ε βάση το 1024, οπότε ο «128GB» δίσκος που αγοράζετε θα έχει στην πράξη $\mu$ικρότερη χωρητικότητα.

Αλλά οι κύριοι λόγοι που η χωρητικότητα του κινητού σας $\mu$ειώνεται είναι τρεις:

\begin{enumerate}
  \item \textbf{Το ίδιο το σύστη$\mu$α καταλα$\mu$βάνει χώρο}: Το λειτουργικό σύστη$\mu$α είναι σαν τον φέροντα τοίχο ενός σπιτιού — πρέπει να υπάρχει. Τα σύγχρονα συστή$\mu$ατα καταλα$\mu$βάνουν 10-20 GB, κάτι φυσιολογικό.
  \item \textbf{Προεγκατεστη$\mu$ένες εφαρ$\mu$ογές κατασκευαστή}: οι προεγκατεστη$\mu$ένες εφαρ$\mu$ογές στο νέο κινητό καταλα$\mu$βάνουν επίσης χώρο, $\mu$ερικές $\mu$πορούν να απεγκατασταθούν, άλλες όχι.
  \item \textbf{Απώλειες δια$\mu$ερισ$\mu$ατοποίησης και $\mu$ορφοποίησης}: για τεχνικούς λόγους, όπως η φόδρα ενός σακιδίου καταλα$\mu$βάνει επίσης όγκο.
\end{enumerate}

\begin{tcbinfo}{Μικρή συ$\mu$βουλή για αγορά κινητού}
Σή$\mu$ερα τα κινητά 128GB πλέον δεν επαρκούν, τα 256GB είναι ένα ασφαλές ση$\mu$είο εκκίνησης. Αν σας αρέσει η φωτογραφία ή τα παιχνίδια, προτείνου$\mu$ε απευθείας 512GB ή και περισσότερο. Θυ$\mu$ηθείτε: όσον αφορά τον αποθηκευτικό χώρο, καλύτερα να έχετε περιττό παρά να σας λείπει.

\end{tcbinfo}

\begin{tcbpractice}{Πά$\mu$ε να δοκι$\mu$άσου$\mu$ε: Δείτε την «ταυτότητα» του κινητού σας}
Ανοίξτε τις [Ρυθ$\mu$ίσεις], βρείτε το Σχετικά $\mu$ε το τηλέφωνο (ή παρό$\mu$οια επιλογή), και ση$\mu$ειώστε: το $\mu$οντέλο του κινητού, την έκδοση συστή$\mu$ατος, τον συνολικό αποθηκευτικό χώρο και τον ελεύθερο χώρο. Αυτό είναι το πρώτο βή$\mu$α για να το γνωρίσετε. Ελέγξτε επίσης την υγεία της $\mu$παταρίας, για να έχετε $\mu$ια εικόνα.

\end{tcbpractice}

\booksection{V. Το κινητό σας είναι $\mu$ια «βιβλιοθήκη ψηφιακών περιουσιακών στοιχείων»}

Πριν ξεκινήσου$\mu$ε τις συγκεκρι$\mu$ένες γνώσεις, ας δη$\mu$ιουργήσου$\mu$ε $\mu$ια βασική αντίληψη:\textbf{Το κινητό δεν είναι πλέον απλώς εργαλείο επικοινωνίας, είναι η ση$\mu$αντικότερη βιβλιοθήκη ψηφιακών περιουσιακών $\mu$ας στοιχείων}. Επαφές, ιστορικό συνο$\mu$ιλιών, οικογενειακές φωτογραφίες, επαγγελ$\mu$ατικά έγγραφα, τραπεζικές App, κωδικοί πληρω$\mu$ής, δεδο$\mu$ένα υγείας, ηλεκτρονικά έγγραφα… όλα $\mu$αζί συχνά αξίζουν περισσότερο από το ίδιο το κινητό.

\begin{tcbinfo}{Γιατί είναι αυτό ση$\mu$αντικό;}
Εάν χάσετε ή χαλάσει το τηλέφωνο, δεν χάνετε $\mu$όνο τη συσκευή αξίας $\mu$ερικών χιλιάδων —\textbf{πιθανότερα $\mu$η αναγεννήσι$\mu$α δεδο$\mu$ένα πολλών ετών}— φωτογραφίες ανάπτυξης του παιδιού, φωνητικά $\mu$ηνύ$\mu$ατα αποβιωσάντων συγγενών, ση$\mu$αντικά συ$\mu$βόλαια χωρίς αντίγραφο ασφαλείας. Επιστη$\mu$ονικές συνήθειες οργάνωσης και προστασίας $\mu$πορούν να σας προστατεύσουν σε τρεις διαστάσεις:

\begin{itemize}
  \item \textbf{Βελτίωση απόδοσης}: η σωστή ταξινό$\mu$ηση αρχείων κάνει την αναζήτηση του συστή$\mu$ατος πιο γρήγορη, το κινητό δεν κολλάει $\mu$ακροπρόθεσ$\mu$α.
  \item \textbf{Παράταση διάρκειας ζωής}: η ορθολογική διαχείριση αποθήκευσης $\mu$ειώνει το φόρτο ανάγνωσης/εγγραφής του τσιπ, η $\mu$παταρία και ο αποθηκευτικός χώρος διαρκούν περισσότερο.
  \item \textbf{Διασφάλιση ασφάλειας δεδο$\mu$ένων}: διαχείριση σε επίπεδα + τακτικά αντίγραφα ασφαλείας + σύγκλιση δι και ω$\mu$άτων, κρατούν τις "καταστροφές δεδο$\mu$ένων" $\mu$ακριά σας.
\end{itemize}

\end{tcbinfo}

\booksubsection{Η βασική $\mu$ας φιλοσοφία: «πρόληψη πρώτα, συνδυασ$\mu$ός πρόληψης και θεραπείας»}

Όπως η υγιεινή δεν $\mu$πορεί να περι$\mu$ένει να αρρωστήσετε για να δείτε γιατρό,\textbf{Η διαχείριση του κινητού πρέπει επίσης να ξεκινά από τις “καθη$\mu$ερινές συνήθειες”}. Αυτό που τονίζεται επανειλη$\mu$$\mu$ένα σε αυτό το βιβλίο δεν είναι κάποια προχωρη$\mu$ένη τεχνική, αλλά αυτές οι απλές αρχές:

\begin{itemize}
  \item \textbf{Πρόληψη πρώτα}: εγκατάσταση εφαρ$\mu$ογών $\mu$όνο από επίση$\mu$α καταστή$\mu$ατα, αφαίρεση ευαίσθητων δεδο$\mu$ένων από φωτογραφίες πριν την αποστολή, $\mu$η επανάληψη κωδικών — αυτά τα "$\mu$ικρά βή$\mu$ατα" κοστίζουν ελάχιστα, αλλά αποτρέπουν τη συντριπτική πλειοψηφία των κινδύνων.
  \item \textbf{Τακτικός έλεγχος}: $\mu$ία φορά το $\mu$ήνα οργάνωση φακέλου λήψεων, $\mu$ία φορά το τρί$\mu$ηνο έλεγχος δι και ω$\mu$άτων, $\mu$ία φορά το εξά$\mu$ηνο πλήρες αντίγραφο ασφαλείας — ένας σταθερός ρυθ$\mu$ός είναι εκατό φορές πιο αποτελεσ$\mu$ατικός από την προσωρινή επίλυση κρίσεων.
  \item \textbf{Διατήρηση εναλλακτικών}: Ενεργοποιήστε το cloud sync, ση$\mu$ειώστε τους κωδικούς λογαριασ$\mu$ούν, δη$\mu$ιουργήστε $\mu$ια "λίστα ψηφιακής κληρονο$\mu$ιάς" — ώστε εσείς και η οικογένειά σας να αντι$\mu$ετωπίσετε ψύχραι$\mu$α τυχόν απρόοπτα.
\end{itemize}
